\section{Limites de Significância e Erros de Decisão}

\begin{frame}{O Limiar de Significância e Sir Ronald Fisher}
    \begin{block}{A Convenção de 5\%}
        Sir Ronald Fisher publicou em 1925 que o limiar de 5\% ($\alpha = 0.05$) é conveniente para julgar se um desvio deve ou não ser considerado significativo.
    \end{block}
    \begin{alertblock}{Cuidado com o Limiar}
        \begin{itemize}
            \item Trata-se de uma convenção histórica, não uma verdade matemática absoluta.
            \item Em conjuntos de dados muito grandes, é recomendável buscar limites muito menores (ex: $0.01$ ou $0.001$) para evitar falsos positivos.
            \item Cuidado com a prática do ``P-hacking'' (manipulação de dados até atingir o $p < 0.05$).
        \end{itemize}
    \end{alertblock}
\end{frame}

\begin{frame}{Erros de Decisão Estatística}
    \begin{block}{Erro do Tipo I (Falso Positivo)}
        Ocorre quando rejeitamos a Hipótese Nula $H_0$, embora ela seja verdadeira. A probabilidade máxima de cometer esse erro é igual ao nível de significância $\alpha$.
    \end{block}
    \begin{block}{Erro do Tipo II (Falso Negativo)}
        Ocorre quando falhamos em rejeitar a Hipótese Nula $H_0$, embora a Hipótese Alternativa $H_1$ seja verdadeira. A probabilidade desse erro é representada por $\beta$.
    \end{block}
    \begin{alertblock}{Poder do Teste}
        Representado por $1 - \beta$, é a capacidade do teste estatístico de detectar um sinal real quando ele realmente existe.
    \end{alertblock}
\end{frame}

\begin{frame}{Matriz de Decisão Estatística}
    \begin{table}
        \centering
        \footnotesize
        \begin{tabular}{lcc}
            \toprule
            \textbf{Realidade} & \textbf{Decisão: Não Rejeitar $H_0$} & \textbf{Decisão: Rejeitar $H_0$} \\
            \midrule
            $H_0$ é Verdadeira & Decisão Correta & \textbf{Erro Tipo I} ($\alpha$ Falso Positivo) \\ \addlinespace
            $H_1$ é Verdadeira & \textbf{Erro Tipo II} ($\beta$ Falso Negativo) & Decisão Correta (Poder: $1-\beta$) \\
            \bottomrule
        \end{tabular}
        \caption{Matriz de decisão e probabilidades de erro}
    \end{table}
    \centering
    \textit{Existe um compromisso entre $\alpha$ e $\beta$: diminuir a tolerância a falsos positivos ($\alpha$) tipicamente reduz o poder do teste, aumentando os falsos negativos ($\beta$).}
\end{frame}

\begin{frame}{Resumo e Principais Aprendizados}
    \begin{block}{Pontos-Chave da Aula}
        \begin{itemize}
            \item \textbf{Modelo Nulo:} O modelo do acaso no qual conseguimos realizar simulações computacionais.
            \item \textbf{TVD:} Estatística de teste robusta para comparar distribuições com múltiplas categorias.
            \item \textbf{Teste de Permutação:} Método de reamostragem ideal para Testes A/B, pois não exige premissas distributivas.
            \item \textbf{P-Valor:} Mede a compatibilidade entre a amostra observada e o modelo nulo.
        \end{itemize}
    \end{block}
    \centering
    \textbf{Dica:} Sempre comece definindo com clareza a Hipótese Nula ($H_0$) e a Estatística de Teste apropriada!
\end{frame}

\begin{frame}{Referências Bibliográficas}
    \begin{block}{Leituras Recomendadas}
        \footnotesize
        \begin{itemize}
            \item \textbf{Learning Data Science} \\
            Sam Lau, Joey Gonzalez, and Deb Nolan. \\
            Capítulo 17: Inference, Prediction, and Generalization. \\
            \url{https://learningds.org/ch/17/inf_pred_gen_HT.html}
            \vspace{0.3cm}
            \item \textbf{Computational and Inferential Thinking} \\
            Ani Adhikari and John DeNero. \\
            Capítulo 11: Testing Hypotheses. \\
            \url{https://inferentialthinking.com/chapters/11/testing-hypotheses/}
        \end{itemize}
    \end{block}
\end{frame}

